% !TEX root = ../tex/1-4.tex
\begin{frame}{1.4 Closing Remarks}
    \begin{block}{この章の目標}
        未知で複雑なデータ分布\(p_{\mathrm{data}}\)を近似する、計算可能なモデル分布\(p_{\phi}\)を学習する。
    \end{block}
    \vspace{0.8em}
    妥当な確率密度を定義するには、正規化定数（分配関数）
    \[
        Z(\phi)=\int \exp(-\mathit{E}_{\phi}(\boldsymbol{x}))\,d\boldsymbol{x}
    \]
    が必要になる。高次元では、この積分が計算できない。
\end{frame}
%========================================
\begin{frame}{計算困難への対処}
    \begin{block}{戦略の違うモデル族}
        \(Z(\phi)\)を直接計算しないために、複数の族が用意されている。
    \end{block}
    \vspace{0.6em}
    この章で見たものは次の五つである。
    \begin{itemize}
        \item Energy-Based Models（EBMs）
        \item Autoregressive Models（ARs）
        \item Variational Autoencoders（VAEs）
        \item Normalizing Flows（NFs）
        \item Generative Adversarial Networks（GANs）
    \end{itemize}
    \vspace{0.4em}
    大きく分けると、密度を計算可能な形で書く Explicit と、\\
    サンプリング手続きだけで分布を定める Implicit になる。
\end{frame}
%========================================
\begin{frame}{拡散モデルの三つの起源}
    \begin{exampleblock}{このモノグラフの焦点}
        古典的な枠組みはどれも重要である。\\
        そのうち三つが、拡散モデルの概念上の出発点になる。
    \end{exampleblock}
    \vspace{0.8em}
    \begin{itemize}
        \item VAEs：変分推論と潜在変数
        \item EBMs：エネルギーとスコア
        \item NFs：可逆な変換
    \end{itemize}
    \vspace{0.6em}
    以降の章では、この三つから拡散モデルがどう展開するかをたどる。
\end{frame}
%========================================
\begin{frame}{Part B：変分の視点}
    \begin{block}{Chapter 2}
        VAEの階層的な潜在変数の構造から、Denoising Diffusion Probabilistic Models（DDPMs）の定式化へ進む。
    \end{block}
    \vspace{0.8em}
    \begin{itemize}
        \item 出発点は変分下界（ELBO）
        \item 潜在変数を一段ではなく、多段にする
        \item その多段の構造が、ノイズを少しずつ除く拡散モデルになる
    \end{itemize}
\end{frame}
%========================================
\begin{frame}{スコアに基づく視点}
    \begin{block}{Chapter 3 から Chapter 4}
        EBMとスコアマッチングから始まり、Noise Conditional Score Networks（NCSN）を経て、より一般の Score SDE に至る。
    \end{block}
    \vspace{0.8em}
    \begin{itemize}
        \item 密度そのものではなく、対数密度の勾配（スコア）を学習する
        \item NCSNは、ノイズの強さごとにスコアを学ぶ
        \item Score SDEは、その手続きを確率微分方程式として書く
    \end{itemize}
\end{frame}
%========================================
\begin{frame}{フローに基づく視点}
    \begin{block}{Chapter 5}
        Normalizing Flowsの原理の上に、生成を連続的な変換として定式化する。\\
        その一般化が Flow Matching である。
    \end{block}
    \vspace{0.8em}
    \begin{itemize}
        \item NFは、単純な分布を可逆変換でデータの分布へ運ぶ
        \item 拡散モデルでは、この変換を連続時間で考える
        \item Flow Matchingは、その変換の速度場を直接学習する
    \end{itemize}
\end{frame}
%========================================
\begin{frame}{この章の位置づけ}
    \begin{exampleblock}{}
        三つの起源を押さえると、拡散モデルの多様な定式化を一つの枠組みとして読める。\\
        それらを貫く原理も、そこから見える。
    \end{exampleblock}
    \vspace{1em}
    \begin{itemize}
        \item 変分 \quad VAE \(\rightarrow\) DDPM（2章）
        \item スコア \quad EBM \(\rightarrow\) NCSN \(\rightarrow\) Score SDE（3--4章）
        \item フロー \quad NF \(\rightarrow\) Flow Matching（5章）
    \end{itemize}
\end{frame}
